\chapter{Conclusions}

In this report, we have explored the fundamental concepts and practical applications of VM and Docker containerization technology.
I began by discussing the architecture of Virtual Machines (VMs) and I then examined the process of creating and managing VMs,
highlighting best practices for optimizing performance and security.
I then briefly discussed Docker, including its components such as Docker Engine, Docker Images, and Docker Containers.
I then delved into the process of creating and managing Docker images, highlighting best practices for building efficient and secure images. \\
Finally, I compared VMs and Docker containers, discussing their respective advantages and disadvantages,
and provided guidance on when to use each technology based on specific use cases and requirements. \\
Overall, this report told that both VMs and Docker containers are powerful tools for managing and deploying applications,
but the performance and efficiency of Docker containers make them a more attractive option for many use cases. \\
Personally, I found Docker even easier to use and manage compared to VMs. The program installation and configuration process was straightforward,
the rapidity of deployment was impressive, and the portability in other systems made it a versatile choice for various applications. \\
Nonetheless, VMs still have their place in certain scenarios, such as when running legacy applications or when a higher level of isolation is required. \\
In conclusion, the choice between VMs and Docker containers ultimately depends on the specific needs and requirements of the application being deployed,
but it is comprehensible that containerization technology is becoming increasingly popular in modern software development and deployment practices.